\section{Algebra \& Number Theory}
\label{sec:prep-algebra}

This section covers the Algebra \& Number Theory track. Items marked \textit{(advanced)} reflect verified exam content and extend beyond the standard undergraduate syllabus.

\subsection{Group Theory and Representations}

\begin{itemize}
  \item Group actions, orbits, stabilizers, conjugacy classes \hfill \cref{algebra:2010:1,algebra:2010:3}
  \item Sylow theorems; simple groups; group of order $150$ not simple \hfill \cref{algebra:2010:11}
  \item $SL_2(\mathbb{C})$ symmetric powers; Clebsch--Gordan decomposition \hfill \cref{algebra:2010:12}
  \item $SL_2(\mathbb{Z})$ generated by $S=\begin{pmatrix}0&-1\\1&0\end{pmatrix}$ and $T=\begin{pmatrix}1&1\\0&1\end{pmatrix}$
  \item $S_4$ irreducible representations; character table construction \hfill \cref{algebra:2010:6}
  \item Molien series for invariant rings
  \item Tensor invariants; $O(V)$-invariants in $V\otimes V$; $\mathbb{C}[\mathrm{End}(V)]^{SL_2\times SL_2}$
  \item Faithful representations; cyclic groups; 1-dimensional reps
  \item Representation theory of $SU(2)$ and $SU(3)$ (Young tableaux basics)
\end{itemize}

\subsection{Galois Theory and Field Extensions}

\begin{itemize}
  \item Field embeddings; $\mathbb{C}\hookrightarrow M_2(\mathbb{R})$, Cayley--Hamilton \hfill \cref{algebra:2010:2}
  \item Galois groups of polynomials: $x^8-5$ (degree 32, quaternion), $x^5-x+1$ ($\mathrm{Gal}\cong S_5$) \hfill \cref{algebra:2010:4}
  \item Solvable groups; solvable Galois groups; radical extensions
  \item Kummer theory; biquadratic polynomials $X^4+aX^2+b$; $(\mathbb{Z}/2)^2$ extensions
  \item Splitting fields; normal extensions; separability
  \item Inertia groups; ramification index and degree; discriminant
  \item Teichm\"uller character: $\mathbb{Z}_p^\times\cong\mathbb{Z}_p\times\mathbb{Z}/(p-1)$
  \item Semidirect products $C_e\rtimes C_f$; absolutely irreducible $\mathbb{F}_p$-representations
  \item Infinite Galois extensions; cyclic subextensions
  \item Normal basis theorem; trace and norm
\end{itemize}

\subsection{Ring and Module Theory}

\begin{itemize}
  \item UFD checks: $\mathbb{Z}[\sqrt6]$, $\mathbb{Z}[(1+\sqrt{-11})/2]$, $\mathbb{C}[x,y]/(x^2+y^2-1)$
  \item Module over a PID; rational canonical form; Jordan canonical form
  \item Free modules; submodules of free modules over PID
  \item Orders in quadratic fields; roots of unity
  \item Weyl algebra; Jacobson density theorem
  \item Integral dependence; going-up theorem; lying-over
  \item Noetherian rings; Hilbert basis theorem; Rees algebras
  \item Five Lemma (homological algebra); Snake Lemma
  \item Nakayama lemma; Cayley--Hamilton theorem (matrix version)
  \item Flat, projective, and injective modules
  \item Morita equivalence; rings of integers in $p$-adic fields
\end{itemize}

\subsection{Commutative Algebra and Algebraic Geometry}

\begin{itemize}
  \item Coordinate rings; $k[x,y]/(y^2-x^2-x^3)$; normalization
  \item Integral closure; normalization of $k[x,y]/(x^2-y^3)$
  \item Hilbert polynomial; Hilbert--Samuel function
  \item Semidualizable extensions $k[X]/(X^p-a)$; nilpotent radical
  \item Transcendence degree; $k((t))/k(t)$ is purely inseparable
  \item Function fields; discrete valuation rings; localizations $k[T]_{(T)}$
  \item Power series rings; $R/(T^n)$ bases in ramified extensions
  \item Algebraic curves; genus; Riemann--Roch (basic)
  \item Separable and inseparable field extensions
\end{itemize}

\subsection{$p$-adic Analysis and Algebraic Number Theory}

\begin{itemize}
  \item $p$-adic valuation; $\mathbb{Q}_p$ completion; spherical completeness
  \item Hensel's lemma; lifting lemmas for roots modulo $p^n$
  \item Contraction maps on $\mathbb{Z}_p$; fixed points of $\phi(x)=x^p+p\cdot g(x)$
  \item $p$-th power in $\mathbb{Z}_p$; Fermat equations over $\mathbb{Z}_7$
  \item $p$-adic logarithm and exponential; $K^\times/K^{\times,n}$ finite
  \item Cyclotomic extensions of $\mathbb{Q}_p$; ramification in $p$-adic local fields
  \item Unramified $p$-adic Galois extensions; faithfully flat descent
  \item Discriminant of cubic fields; integral basis
  \item Mordell curves $x^2+13=y^3$ (class number 2)
  \item Pell equation $X^2-82Y^2=\pm2$; Hasse--Minkowski
  \item Dirichlet unit theorem; regulator; fundamental units
  \item Mason--Stothers theorem (abc for polynomials)
  \item Quadratic reciprocity; Gauss sums; cyclotomic fields $\mathbb{Q}(\zeta_p)$
\end{itemize}

\subsection{Linear and Multilinear Algebra}

\begin{itemize}
  \item Nilpotent operators; Jordan form; common eigenvector for commuting matrices \hfill \cref{algebra:2010:1}
  \item Cayley--Hamilton theorem (polynomial and matrix forms)
  \item Tensor product; symmetric and exterior powers; $\Lambda^r V^*$
  \item Cartan--Dieudonn\'e theorem ($O(n)$ generated by reflections)
  \item Positive semidefinite matrices $t^{a_i+a_j}$; determinant inequalities
  \item Skew-symmetric matrices; $e^A$ orthogonal
  \item Normal operators; unitary diagonalization
  \item Trace conditions $\mathrm{tr}(A^n)\in F\Rightarrow\mathrm{tr}(A)\in F$
  \item Determinants; Cauchy determinant formula
\end{itemize}

\subsection{Combinatorics and Elementary Number Theory}

\begin{itemize}
  \item Generating functions; Beatty sequences; partition generating functions \hfill \cref{algebra:2010:5,algebra:2010:7}
  \item Catalan numbers; Ballot theorem; Cycle Lemma
  \item Mason--Stothers theorem; Mason's equation $a+b=c$ \hfill \cref{algebra:2010:8}
  \item Chevalley--Warning theorem ($p$-adic methods)
  \item Fibonacci sequences modulo $p$; quadratic norm maps
  \item Mersenne and Fermat primes; Goldbach coprimality
  \item Congruences; Chinese remainder theorem; primitive roots
  \item Trigonometric products $\prod\cos(n\pi/p)$ \hfill \cref{algebra:2010:9}
  \item Finite fields $\mathbb{F}_p$; irreducible polynomials over $\mathbb{F}_2$
  \item Radon transform; union of proper subspaces
  \item Probabilistic method; conjugacy class bound $c(G)\leq 5|G|/8$ \hfill \cref{algebra:2010:3}
\end{itemize}

\section{Analysis \& Differential Equations}
\label{sec:prep-analysis}

This section covers both the official syllabus and graduate-level topics that recur in actual exams. Items marked \textit{(advanced)} go beyond the undergraduate syllabus.

\subsection{Real Analysis and Measure Theory}

\begin{itemize}
  \item Uniform convergence; Dini's theorem; Arzel\`a--Ascoli
  \item Lebesgue $L^p$ spaces; H\"older and Minkowski inequalities
  \item Fatou's lemma; dominated and monotone convergence theorems
  \item Riesz--Markov representation theorem
  \item Urysohn lemma; Tietze extension theorem
  \item Layer-cake representation
  \item Weak $L^p$ spaces; tail estimates
  \item Product measures; Fubini--Tonelli
  \item Signed measures; Radon--Nikodym theorem
  \item Lebesgue differentiation theorem; approximation of identity
  \item Borel--Cantelli lemma; limsup of events
\end{itemize}

\subsection{Fourier Analysis}

\begin{itemize}
  \item Fourier series: $L^2$ convergence; Parseval identity
  \item Dirichlet and Fej\'er kernels; pointwise convergence; Gibbs phenomenon
  \item Weyl equidistribution theorem
  \item Fourier uncertainty principle
  \item Fourier restriction and Strichartz estimates
  \item Paley--Wiener theorem
  \item Hardy space $H^2$; boundary values of analytic functions
  \item Orthonormal bases in Hilbert spaces
  \item Fej\'er and Ces\`aro summability
\end{itemize}

\subsection{Complex Analysis}

\begin{itemize}
  \item Cauchy integral theorem and formula; Morera's theorem
  \item Power series; radius of convergence; analytic continuation
  \item Contour integration: keyhole, annulus-to-strip, wedge contours \hfill analysis:2011:1,analysis:2011:3
  \item Residue theorem; argument principle; Rouch\'e's theorem
  \item Conformal mapping; Riemann mapping theorem
  \item Schwarz--Pick lemma; Jenson\'s formula
  \item Winding number; homotopy invariance
  \item Weierstrass factorization; Hadamard factorization
  \item Maximum modulus principle; Phragm\'en--Lindel\"of method
  \item Hadamard three-circles theorem
  \item Finite Blaschke products; product decomposition
  \item Elliptic functions; elliptic function fields
\end{itemize}

\subsection{Functional Analysis}

\begin{itemize}
  \item Hilbert spaces; orthonormal systems; Riesz representation theorem
  \item Dual spaces; Hahn--Banach theorem (analytic and geometric forms)
  \item Baire category theorem; open mapping and closed graph theorems
  \item Uniform boundedness principle
  \item Compact operators; spectral theory for self-adjoint compact operators \hfill analysis:2011:6,analysis:2011:7
  \item Volterra operator compactness
  \item Fredholm alternative
  \item Sobolev spaces $H^s$; Rellich--Kondrachov embedding \hfill \cref{analysis:2011:6}
  \item BMO; John--Nirenberg inequality; doubling measures
\end{itemize}

\subsection{Partial Differential Equations}

\begin{itemize}
  \item Laplace equation: Poisson integral; Dirichlet and Neumann problems
  \item Heat equation: fundamental solution; maximum principle
  \item Wave equation: d'Alembert formula; energy method
  \item Method of characteristics for first-order PDE
  \item Semilinear elliptic equations; sub- and super-solutions; contraction mapping
  \item Boltzmann equation basics; mild solutions
  \item Caccioppoli inequality
  \item Interior Schauder estimates
  \item Picard iteration for holomorphic maps to high-genus curves
\end{itemize}

\subsection{Harmonic Analysis and Potential Theory}

\begin{itemize}
  \item Poisson kernel; mean value property for harmonic functions
  \item Harnack inequality and Harnack convergence theorem
  \item Subharmonic and superharmonic functions
  \item Green's function; balayage
  \item Newtonian and Riesz potentials
  \item B\^ocher's theorem
  \item Newton potential; Poisson equation $\Delta u = f$
\end{itemize}

\subsection{Calculus of Variations and Convex Analysis}

\begin{itemize}
  \item Direct method; lower semicontinuity
  \item Euler--Lagrange equations; natural boundary conditions
  \item Convexity; Jensen's inequality (function convexity)
  \item Legendre transform; dual problems
  \item Fundamental lemma of calculus of variations
  \item Cartan's uniqueness theorem
  \item Dirichlet principle
\end{itemize}

\subsection{ODEs and Qualitative Analysis}

\begin{itemize}
  \item Existence and uniqueness; Picard--Lind\"of theorem
  \item Linear ODEs; fundamental systems; Wronskian \hfill \cref{analysis:2011:2}
  \item Matrix exponentials; $e^{tA}$
  \item Baker--Campbell--Hausdorff formula
  \item Trotter product formula
  \item Phase portraits; stability of equilibria; Floquet theory
  \item Lyapunov stability; Gronwall inequality
  \item Hill's equation; periodic solutions \hfill \cref{analysis:2011:8}
  \item Hamiltonian systems; action-angle variables
  \item Euler--Cauchy equation
\end{itemize}

\subsection{Differential and Integral Calculus on Euclidean Space}

\begin{itemize}
  \item Partial derivatives; Clairaut's theorem; tangent plane
  \item Chain rule; gradient and Hessian; Taylor's theorem in several variables
  \item Inverse and implicit function theorems; rank theorem
  \item Lagrange multipliers; constrained optimization
  \item Double and triple integrals; Fubini's theorem
  \item Change of variables formula; Jacobian
  \item Surface integrals; area of parametric surfaces
  \item Line integrals; conservative vector fields; path independence
  \item Green's theorem; divergence theorem; Stokes' theorem
  \item Improper integrals; convergence criteria
\end{itemize}

\section{Geometry \& Topology}
\label{sec:prep-geometry}

\textbf{Note:} The official syllabus lists only high-school/undergraduate geometry. The actual examination is at the graduate level. All topics below are drawn from verified past examinations.

\subsection{Algebraic Topology}

\begin{itemize}
  \item Fundamental group and covering spaces
  \item Singular homology and cohomology
  \item de Rham cohomology; Mayer--Vietoris sequence
  \item Suspension homology: $\widetilde{H}_k(SX)\cong\widetilde{H}_{k-1}(X)$
  \item Brouwer fixed-point theorem and Borsuk--Ulam theorem
  \item Jordan--Brouwer separation theorem
  \item Degree of a map $S^n\to S^n$; Hopf degree theorem
  \item Hopf fibration $S^1\hookrightarrow S^3\to S^2$; $\pi_3(S^2)\neq0$ \hfill \cref{geometry:2012:1}
  \item Adams--Hopf invariant one
  \item Essential maps $M\to S^n$; $M\to T^n$ (homotopy vanishing)
  \item Linking number; linking of circles in $\mathbb{R}^3$
  \item Cellular homology; CW complexes
  \item Homology of $m$-skeleton of $n$-simplex
\end{itemize}

\subsection{Differential Topology}

\begin{itemize}
  \item Poincar\'e--Hopf index theorem
  \item Pontryagin--Thom construction
  \item Characteristic classes: Stiefel--Whitney, Pontryagin, Euler
  \item Cobordism theory; $\mathbb{RP}^k$ as boundaries
  \item Tangent bundle triviality and Euler characteristic
  \item Immersion obstructions; Pontryagin classes
  \item $\mathbb{CP}^2\not\hookrightarrow\mathbb{R}^6$ (Pontryagin--Thom obstruction)
  \item $\mathbb{CP}^2$ orientation
  \item H-space; Adams--Hopfan obstruction
  \item Reduced suspension; adjunction spaces
\end{itemize}

\subsection{Differential Geometry}

\begin{itemize}
  \item Riemannian comparison geometry: Rauch, Bishop--Grovov volume comparison
  \item Myers's diameter theorem; Cheng's eigenvalue bound
  \item Alexandrov spaces with curvature bounded below
  \item Constant mean curvature spheres; Liebmann's theorem (Alexandrov)
  \item Gauss--Bonnet theorem for surfaces \hfill \cref{geometry:2012:12}
  \item Chern--Weil theory; trace of curvature polynomials is closed
  \item Normal coordinates; curvature expansions
  \item Levi-Civita connection; Koszul formula (uniqueness) \hfill \cref{geometry:2012:9}
  \item Riemann structure equations; curvature 2-forms
  \item Minkowski formulas for surfaces in $\mathbb{R}^3$ \hfill \cref{geometry:2012:5}
  \item Bochner formula; sphere eigenfunctions
  \item Frankel's theorem: geodesics on positively curved surfaces intersect
  \item Mean curvature in $\mathbb{H}^3$ (conformal change formula)
  \item Simons inequality for minimal surfaces in $S^{n+p}$
  \item Bernstein's theorem (minimal graphs in $\mathbb{R}^3$)
  \item Huisken shrinkers; $S^n(\sqrt{2n})$
  \item Contact structures; $\mathbb{RP}^{2n+1}$ contact structure
  \item Berger spheres; $S^3$ metrics $g_\varepsilon$
  \item Conformal geometry; curvature equation $\Delta u-K+\widetilde{K}e^{2u}=0$
  \item Keller--Osserman condition; no entire solutions of $\Delta u=e^u$
  \item Wirtinger inequality (isoperimetric inequality via Wirtinger)
  \item Flow-box coordinates; Frobenius theorem (integrability)
  \item Bertrand curves; tube surfaces
\end{itemize}

\subsection{Lie Groups and Transformation Groups}

\begin{itemize}
  \item Lie group topology: $\pi_1(SO(3))$, $SU(n)$, $SO(n)$
  \item $SO(3)\cong\mathbb{RP}^3$; $\pi_1(SO(n))\cong\mathbb{Z}/2$ for $n\geq3$
  \item Unitary group homotopy: $\pi_2(U(n))=0$, $\pi_3(U(n))\cong\mathbb{Z}$
  \item $\det_*: \pi_1(U(n))\cong\pi_1(S^1)\cong\mathbb{Z}$
  \item Isometries on compact metric spaces; surjectivity
  \item Transformation groups; fixed-point-free involutions on $\mathbb{CP}^n$
\end{itemize}

\subsection{Global Analysis and Geometric PDE}

\begin{itemize}
  \item Hodge theory; harmonic forms; $H^1(M;\mathbb{R})=0$ under positive Ricci
  \item Signature of 4-manifolds ($S^4$, $\mathbb{CP}^2$, $S^2\times S^2$)
  \item First eigenvalue of minimal hypersurface in $S^{n+1}$; $\lambda_1\geq n/2$
  \item Conformal representation of surfaces
  \item Elliptic PDE on manifolds; sub- and super-solution method
  \item Eigenvalue variation under metric deformation
\end{itemize}

\section{Probability \& Statistics}
\label{sec:prep-probability}

This section covers both the official syllabus and advanced topics from past exams. Items marked \textit{(advanced)} go beyond undergraduate-level probability.

\subsection{Foundations of Probability}

\begin{itemize}
  \item Sample space, events, classical probability, geometric probability, probability axioms, conditional probability, Bayes' theorem, independence
  \item Borel--Cantelli lemma
  \item Convergence theorems: convergence in probability, in distribution, in $L^p$, almost sure convergence, SLLN, CLT, law of iterated logarithm (LIL), Berry--Esseen
  \item Concentration inequalities: Hoeffding, Khinchin, Rademacher--Menshov
\end{itemize}

\subsection{Random Variables and Distributions}

\begin{itemize}
  \item Discrete: binomial, Poisson, hypergeometric, negative binomial, geometric
  \item Continuous: uniform, exponential, normal (Gaussian), gamma, beta, Dirichlet
  \item Distribution functions, pmf/pdf, transformation of random variables
  \item Multivariate distributions; joint and marginal; conditional distributions; convolution
  \item Covariance, correlation, independence; MGF, CF, moments
  \item Order statistics; records; exponential minima; Stein's lemma
  \item Multivariate normal; chi-squared, $t$, $F$ distributions
\end{itemize}

\subsection{Mathematical Expectation}

\begin{itemize}
  \item Expectation, variance, standard deviation
  \item Conditional expectation; tower property
  \item Wald's equation; recursive equation; generating function method
  \item Martingale; optional stopping; martingale convergence
  \item Random walk; gambler's ruin; hitting probabilities
  \item Queueing theory (M/M/1, birth--death); renewal theory; overlapping patterns
  \item Coupling; total variation distance
  \item Foster--Lyapunov drift criterion
\end{itemize}

\subsection{Stochastic Processes}

\begin{itemize}
  \item Markov chains: state space, transition matrix, Chapman--Kolmogorov, recurrent/transient classification, stationary distribution, limit distribution, reversibility, ergodicity
  \item Branching process; extinction probability
  \item Poisson process: construction, interarrival times, superposition, thinning
  \item Brownian motion (Wiener process): martingale properties, hitting times, reflection principle
  \item Urn models: Polya's urn, reinforcement
  \item Random graphs: Erd\H{o}s--R\'enyi $G(n,p)$, first and second moment method, phase transition, connectivity threshold, largest clique
\end{itemize}

\subsection{Statistical Estimation and Decision Theory}

\begin{itemize}
  \item MLE; score function; Fisher information; asymptotic efficiency
  \item UMVUE; Cram\'er--Rao lower bound
  \item James--Stein shrinkage; admissibility
  \item Decision theory: loss, risk, Bayes risk, minimax risk; optimal testing (Neyman--Pearson); confidence intervals; minimax intervals
  \item GLM: link function, logistic regression, Poisson regression
  \item Lasso; $L^1$ regularization; KKT conditions
  \item Kernel density estimation (Epanechnikov)
  \item Nonparametric tests: Wilcoxon signed-rank, Wilcoxon--Mann--Whitney, permutation tests, Fisher exact test
\end{itemize}

\subsection{Causal Inference and Experimental Design}

\begin{itemize}
  \item Neyman--Rubin causal model; potential outcomes; SUTVA; ATT vs ATE
  \item Fisher exact test for $H_0$
  \item Randomized experiment; randomization inference; re-randomization
  \item IPW; propensity score; covariate adjustment; ANCOVA
\end{itemize}

\section{Applied \& Computational Mathematics}
\label{sec:prep-applied}

This section covers both the official syllabus and computationally-focused topics from past exams. Items marked \textit{(advanced)} reflect verified exam content.

\subsection{Calculus (syllabus)}

\begin{itemize}
  \item Functions, limits, continuity
  \item Derivatives: tangent lines, optimization, related rates
  \item Integrals: indefinite, definite, improper; FTC; techniques (substitution, parts, partial fractions, trig substitution)
  \item Sequences and series; convergence tests; power series; Taylor series
  \item Multivariable: partial derivatives, gradient, divergence, curl, Jacobian
  \item Multiple integrals; line and surface integrals
  \item Green, Gauss, Stokes theorems
  \item ODEs: separable, linear, exact, integrating factor
  \item Higher-order linear ODEs with constant coefficients
  \item Systems of first-order ODEs; phase portraits
  \item PDEs: classification (elliptic, parabolic, hyperbolic); canonical equations (heat, wave, Laplace)
\end{itemize}

\subsection{Linear Algebra (syllabus)}

\begin{itemize}
  \item Determinants; cofactor expansion
  \item Matrix operations; inverse matrices; Sherman--Morrison rank-one update
  \item Gaussian elimination; LU; RREF
  \item Vector spaces: subspaces, basis, dimension, coordinates
  \item Linear independence; rank--nullity
  \item Eigenvalues and eigenvectors; characteristic polynomial; Cayley--Hamilton
  \item Spectral theorem for symmetric matrices
  \item Inner product; Gram--Schmidt; QR decomposition; least squares
  \item Quadratic forms; Gershgorin disc theorem
\end{itemize}

\subsection{Probability (syllabus)}

\begin{itemize}
  \item Sample spaces, events, probability axioms; conditional probability; Bayes
  \item Random variables: discrete and continuous distributions
  \item Joint distributions; marginal and conditional
  \item Expectation, variance, covariance, correlation; MGF and CF
  \item LLN; CLT; conditional expectation
\end{itemize}

\subsection{Operations Research (syllabus)}

\begin{itemize}
  \item LP: formulation, geometry, simplex, two-phase, big-M, duality; sensitivity analysis
  \item Integer LP; branch-and-bound
  \item Transportation and assignment; Hungarian algorithm
  \item Network flows: Dijkstra, Floyd--Warshall, Kruskal, Prim, Ford--Fulkerson
  \item CPM/PERT; decision analysis: decision trees, expected value, risk
  \item Game theory: zero-sum, Nash equilibrium, mixed strategies
\end{itemize}

\subsection{Finite Difference Methods for PDEs} \textit{(advanced)}

\begin{itemize}
  \item Consistency, stability, convergence; Lax equivalence theorem
  \item Von Neumann stability analysis; CFL condition
  \item Explicit schemes: FTCS for heat equation; stability constraints
  \item Upwind scheme; Beam--Warming; Lax--Wendroff; leapfrog scheme
  \item Implicit schemes: BTCS; Crank--Nicolson ($\theta$-scheme); unconditional stability
  \item Richardson extrapolation
  \item Finite difference for wave equation and Laplace/Poisson
\end{itemize}

\subsection{Finite Element Methods (FEM)} \textit{(advanced)}

\begin{itemize}
  \item Weak (variational) formulation of BVPs; Ritz--Galerkin
  \item CG and DG methods; piecewise polynomial spaces ($P_1$, $P_2$)
  \item A priori error estimates; convergence rates in $H^1$ and $L^2$
  \item Superconvergence at Gaussian points
  \item FEM for elliptic BVPs: Poisson equation
\end{itemize}

\subsection{Numerical Linear Algebra} \textit{(advanced)}

\begin{itemize}
  \item SVD; low-rank approximation via SVD; Eckart--Young
  \item QR: Gram--Schmidt (classical and modified), Householder, Givens
  \item Vandermonde least-squares approximation
  \item Eigenvalue algorithms: power method; inverse iteration; shifted power method; QR algorithm
  \item Arnoldi process; Krylov subspaces; Lanczos
  \item CG; GMRES; Richardson; SOR; Jacobi; Gauss--Seidel
  \item Spectral radius of tridiagonal (Toeplitz) matrices; Gershgorin discs
\end{itemize}

\subsection{Numerical Methods for ODEs} \textit{(advanced)}

\begin{itemize}
  \item Forward Euler, backward Euler, improved Euler, RK4
  \item Consistency, order, convergence; absolute stability regions
  \item BDF; BDF2 zero-stability and order; SSP-RK; IMEX RK
  \item Multistep: Adams--Bashforth, Adams--Moulton
  \item Stiff ODEs; A-stability, L-stability
  \item Banach fixed-point for ODE convergence
  \item Secant method (superlinear convergence); Steffensen acceleration; bisection; Newton's method
\end{itemize}

\subsection{Optimization and Convex Analysis} \textit{(advanced)}

\begin{itemize}
  \item KKT conditions; complementary slackness; dual feasibility
  \item Lagrangian duality; weak and strong duality; Slater's condition
  \item Lasso regression; $L^1$ regularization; compressed sensing
  \item Convex sets and functions; co-coercivity; proximal operators; FBP; Douglas--Rachford
  \item Subgradient and subdifferential calculus; Frank--Wolfe basics
\end{itemize}

\subsection{Calculus of Variations and Euler--Lagrange} \textit{(advanced)}

\begin{itemize}
  \item Euler--Lagrange equation; necessary and sufficient conditions; natural boundary conditions
  \item Isoperimetric problems; isoperimetric inequality
  \item Legendre condition; Jacobi conjugate points
  \item Noether's theorem; direct method; Ritz method
\end{itemize}

\subsection{Orthogonal Polynomials and Spectral Methods} \textit{(advanced)}

\begin{itemize}
  \item Legendre polynomials: orthogonality, Rodrigues formula, recurrence
  \item Chebyshev ($T_n$, $U_n$): orthogonality with weight $(1-x^2)^{-1/2}$; minimax property
  \item Gaussian quadrature: Gauss--Legendre, Gauss--Chebyshev; error estimates
  \item Clenshaw--Curtis; Lobatto; Radau
  \item Chebyshev interpolation; spectral collocation
  \item Moving least squares (MLS)
  \item Fourier spectral methods; spectral accuracy
  \item Crank--Nicolson spectral for diffusion
\end{itemize}

\subsection{Dynamical Systems and Biological Modeling} \textit{(advanced)}

\begin{itemize}
  \item Phase plane; fixed points; linearization via Jacobian; stability; Lyapunov stability
  \item Bifurcations: saddle-node, transcritical, pitchfork, Hopf
  \item FitzHugh--Nagumo neuronal model; excitability
  \item Michaelis--Menten enzyme kinetics; steady-state approximation; Lineweaver--Burk plot
  \item Singular perturbation; boundary layers; matched asymptotic expansions
  \item WKB for Schr\"odinger equations
  \item Hamilton--Jacobi equations; shock formation; Rankine--Hugoniot conditions
\end{itemize}

\subsection{Graph Theory (Advanced)}

\begin{itemize}
  \item Eulerian and Hamiltonian graphs; Dirac's theorem ($n\geq3$, $\deg(v)\geq n/2$); Ore's theorem
  \item Bondy--Chvatal closure
  \item Hall's marriage theorem; Tutte's theorem
  \item Planar graphs; Kuratowski ($K_5$, $K_{3,3}$)
  \item Graph coloring; chromatic polynomial
\end{itemize}

\subsection{Number Theory}

\begin{itemize}
  \item Divisibility; Euclidean algorithm; prime numbers
  \item Congruences; Chinese remainder theorem
  \item Fermat's little theorem; Euler's theorem; $\varphi(n)$
  \item RSA encryption; primitive roots; discrete logarithms
  \item Quadratic residues; Legendre and Jacobi symbols
  \item Eisenstein's criterion
\end{itemize}

\subsection{Learning Theory and VC Dimension} \textit{(advanced)}

\begin{itemize}
  \item PAC learning framework; sample complexity
  \item VC dimension: definition, growth function, Sauer's lemma
  \item VC dimension of interval classifiers in $\mathbb{R}$ equals 2
  \item Rademacher complexity basics
\end{itemize}

\section{Mathematical Physics}
\label{sec:prep-physics}

The Mathematical Physics track was first introduced in 2022. The official syllabus is broad, but actual exams focus on a core of five areas.

\subsection{Classical Mechanics}

\begin{itemize}
  \item Lagrangian and Hamiltonian formalisms; polar coordinates; effective 1D reduction
  \item Central potential $V(r)=\alpha r^k$; conserved angular momentum and energy
  \item Circular orbits; effective potential; orbital stability; small oscillations
  \item Constrained motion; Lagrange multipliers; rotating hoop dynamics
  \item Noether's theorem; symmetries and conserved currents \hfill \cref{physics:2022:2}
  \item Aharonov--Bohm magnetic field; gauge-invariant Lagrangian
\end{itemize}

\subsection{Quantum Mechanics}

\begin{itemize}
  \item Harmonic oscillator: ladder operators $\hat a$, $\hat a^\dagger$; $su(2)$ algebra; number states \hfill \cref{physics:2022:1}
  \item Coherent states; displacement operator; transition probabilities
  \item Time-dependent perturbation; Heisenberg picture; equations of motion
  \item Perturbation theory: degenerate and non-degenerate; first-order corrections
  \item Anisotropic harmonic oscillator; exact diagonalization; perturbation correction
  \item Shifted harmonic oscillator; displacement operator
  \item Quantum scattering: step potential; barrier potential; reflection and transmission coefficients
  \item Time-reversal symmetry; anti-unitary operators; Wigner theorem
\end{itemize}

\subsection{Electromagnetism}

\begin{itemize}
  \item Maxwell equations in vacuum and media; constitutive relations
  \item Wave propagation in conductors; complex conductivity; skin depth
  \item Skin effect; wave reflection at conducting half-space
  \item Scalar and vector potentials; Lorentz gauge
  \item Retarded Green's function; retarded potentials
  \item Li\'enard--Wiechert potentials for a moving point charge
  \item Gaussian wave packet; paraxial approximation
\end{itemize}

\subsection{Statistical Mechanics and General Relativity}

\textbf{Statistical Mechanics:}
\begin{itemize}
  \item Ising model: mean-field; Landau free energy; critical exponents $\alpha,\beta,\gamma$
  \item Ising model on triangular lattice: exact partition function (Onsager); specific heat; susceptibility
  \item Free massless bosons; heat capacity; Stefan--Boltzmann
  \item Fermi--Dirac statistics; Sommerfeld expansion; degenerate Fermi gas
  \item Partition function; thermodynamic potentials; magnetization fluctuations
\end{itemize}

\textbf{General Relativity:}
\begin{itemize}
  \item Schwarzschild metric; geodesics; effective potential
  \item Schwarzschild--de Sitter metric; vacuum Einstein equations $R_{\mu\nu}=kg_{\mu\nu}$
  \item Killing vectors; conserved quantities along geodesics \hfill \cref{physics:2022:2}
  \item de Sitter space: static coordinates; geodesic equation; Penrose diagram
  \item Gravitational waves; binary black hole; quadrupole formula; orbital decay
  \item Null hypersurfaces; Eddington--Finkelstein coordinates \hfill \cref{physics:2022:5}
\end{itemize}

\subsection{Quantum Field Theory}

\begin{itemize}
  \item $\phi^3$ theory: Feynman diagrams; one-loop self-energy; dimensional regularization
  \item Yukawa theory: scalar and fermion self-energy at one loop; counterterms
  \item $\phi^4$ theory: loop expansion; scalar self-energy diagrams; $O(N)$ scalar theory; large-$N$ expansion
  \item Dimensional regularization; $\overline{\mathrm{MS}}$ scheme; renormalization conditions
  \item Scale invariance; conformal invariance; trace of the energy-momentum tensor
  \item Trace anomaly; conformal scalar field on curved spacetime; Weyl invariance
\end{itemize}
